\documentclass{article}
\usepackage[preprint]{neurips_2025}

\usepackage[utf8]{inputenc}
\usepackage[T1]{fontenc}
\usepackage{hyperref}
\usepackage{url}
\usepackage{booktabs}
\usepackage{amsfonts, amsmath, amssymb}
\usepackage{nicefrac}
\usepackage{microtype}
\usepackage{xcolor}
\usepackage{graphicx}
\usepackage{subcaption}
\usepackage{natbib}

\title{Ternary Sufficient-Statistic Import for State-Space Language Model Training}

\author{
  Julio Jose Lena \\
  Independent Researcher \\
  \texttt{juliolena25@gmail.com} \\
  \url{https://orcid.org/0009-0001-0817-6591}
}

\begin{document}

\maketitle

\begin{abstract}
This paper introduces Ternary Sufficient-Statistic Import (TSSI), a training
object for state-space language models with ternary weights. The motivating
question is deliberately narrow: before iterative optimization begins, can a
corpus be compiled into model-native $\{-1,0,+1\}$ commitments that already
answer part of the cross-entropy descent problem? TSSI does not import text,
examples, retrieval memory, or a teacher model. It streams the corpus once,
accumulates original cross-entropy directional moments in a fixed SSM feature
map, realizes stable directions as sparse ternary sign/zero commitments, and
then leaves the remaining error to ordinary residual training on the same
next-token objective. Initial gates on byte-level v12 shards show that
corpus-derived ternary commitments beat density-matched random and
sign-shuffled controls, including over live Vulkan SSM hidden states and
record-disjoint held-out trace rotations. The result supports TSSI as a
training-theoretic mechanism: useful corpus information can be converted into a
ternary SSM artifact without becoming raw memorization or a replacement
objective.
\end{abstract}

\section{Introduction}

The usual framing of language-model training treats the corpus as something
that must be visited many times by an optimizer until useful structure has been
written into the parameters. For continuous weights this is natural: the model
must discover both a direction and a magnitude in a high-dimensional real
space. For a ternary state-space model, however, the first useful question is
smaller and more discrete. A coordinate often only needs an initial answer to
one of three possibilities: negative, zero, or positive. If the corpus already
contains enough directional evidence for some of those answers, then forcing
all of that evidence to pass through ordinary training steps is not a law of
learning; it is a design choice.

The phrase ``dataset import'' is easy to misread. A literal import of a
dataset into a model is not the object studied here. Storing examples, building
a lookup table, or hiding retrieval behind the parameter interface would evade
the training problem rather than improve it. The admissible object is narrower:
a sufficient statistic for a particular training objective. In this paper the
objective remains next-token cross-entropy, and the imported artifact is only a
set of sparse ternary commitments in model coordinates. Anything not supported
by stable corpus evidence is left as zero, so the ordinary residual training
process still owns the unsolved part of the model.

This gives a different way to ask what a corpus can do before training. Instead
of asking whether text can be compressed into weights, TSSI asks whether
cross-entropy moments can be compiled into the same ternary decision surface
that training would otherwise have to reach by repeated updates. The compiler
streams contexts and targets, accumulates directional moments under the
original loss, applies an abstention threshold, and writes only
$\{-1,0,+1\}$ decisions. The zero state is important: it is not failure, but an
explicit refusal to pretend that weak evidence is already knowledge.

The empirical claim is correspondingly scoped. On byte-level v12 shards, a
bigram feature gate, residual cross-entropy pilots, a deterministic recurrent
\texttt{ssm\_state} proxy, live Vulkan SSM hidden-state traces, and
record-disjoint held-out rotations all show the same qualitative behavior:
corpus-derived ternary commitments carry real predictive signal beyond
density-matched random or sign-shuffled controls. The live SSM result matters
most because it places the statistic in actual state-space coordinates rather
than in a purely symbolic proxy. The record-disjoint rotations matter because
they reduce the simplest memorization explanation: the held-out records change,
yet the all-layer imported commitments remain better than matched random
commitments.

The contribution is therefore not a new dataset, not a teacher-distillation
method, and not an associative memory. It is a training-theoretic object for
ternary SSMs: a corpus-to-commitment map that preserves the original
cross-entropy objective and initializes part of the ternary search problem
before residual optimization. The rest of the paper defines this object,
separates it from nearby prior art, states the validation protocol, and reports
the current gates.

\section{Background and Related Work}

\paragraph{Dataset distillation and condensation.}
Dataset distillation asks for a small synthetic training set that can replace a
larger corpus for downstream optimization \citep{wang2018datasetdistillation}.
Text-level variants such as DiLM move this idea into language by generating
synthetic text that trains a model efficiently \citep{dilm2024}. TSSI is not a
synthetic dataset. It has no generated examples and no second corpus that the
optimizer later consumes. The output is a model-coordinate artifact: a sparse
ternary commitment pattern derived from the original cross-entropy moments.
The distinction matters because a distilled dataset still delegates learning to
future optimization over examples, while TSSI answers a bounded subset of the
ternary parameter decisions directly.

\paragraph{Data-dependent initialization.}
Another nearby line uses data to choose a better initialization for continuous
training. Classical data-dependent initializations match early activation or
gradient scales \citep{krahenbuhl2015data}; GradInit tunes initialization to
improve early optimization dynamics \citep{zhang2021gradinit}; AutoInit searches
for initialization rules that stabilize neural networks automatically
\citep{bingham2021autoinit}. These methods remain continuous in spirit. Their
goal is to place the optimizer in a numerically favorable region. TSSI instead
uses data to make discrete ternary commitments. The imported state is not a
scale, variance, or conditioning adjustment; it is a partial answer to the
model's sign/zero decision problem under the original loss.

\paragraph{Ternary and low-bit language models.}
Low-bit language-model work shows that large networks can preserve useful
behavior under severe quantization \citep{wang2023bitnet, ma2024era}. Earlier
ternary-weight methods also study how to train or quantize networks with
$\{-1,0,+1\}$ weights \citep{li2016ternary}. Those results make the present
question plausible, but they do not define a corpus-to-ternary import object.
The usual quantization problem starts with trained or trainable real-valued
parameters and asks how to represent them cheaply. TSSI starts one step earlier:
given a corpus and a fixed SSM feature map, which ternary coordinates already
have enough evidence to be committed before residual training continues?

\paragraph{State-space coordinates.}
The coordinate system is also part of the claim. Modern SSMs such as Mamba
derive much of their sequence behavior from recurrent hidden states rather than
only from local token features \citep{gu2024mamba}. A corpus statistic that
passes in a bag-of-symbols proxy may still fail when placed into live recurrent
state coordinates. For that reason, the strongest current TSSI evidence is not
the byte/bigram gate by itself, but the live Vulkan SSM hidden-state trace and
the record-disjoint rotations over that trace. The artifact is intended to be
model-native, not merely corpus-native.

\paragraph{Memory and teacher boundaries.}
TSSI should also be separated from retrieval, associative memory, and
teacher-based distillation. A retrieval system keeps examples available at
inference time. A teacher system transfers behavior from a stronger external
model. TSSI does neither. The compiler sees only corpus tokens and fixed SSM
feature coordinates, writes ternary sign/zero commitments, and then disappears.
After import, the remaining burden is still optimized by the same original
cross-entropy objective. This residual contract is the main semantic guardrail:
if an import path cannot preserve it, it is not TSSI as defined in this paper.

\section{Ternary Sufficient-Statistic Import}

TSSI is a corpus compiler for a fixed model-coordinate view of the training
objective. It should not be read as an attempt to solve the model in closed
form. The compiler only asks whether the corpus gives enough first-order
evidence to commit some ternary coordinates before the residual optimizer takes
over.

\subsection{Corpus Moment Accumulation}

Let $D = \{(c_t, y_t)\}_{t=1}^T$ denote a stream of next-token contexts and
targets. Let $\psi(c_t) \in \mathbb{R}^m$ be a fixed feature map in the
coordinate system where the import will be written. Depending on the gate,
$\psi$ can be a byte bigram feature, a deterministic recurrent proxy, or a live
SSM hidden-state summary. Let $p_0(\cdot \mid c_t)$ be the reference model
distribution before import. The cross-entropy gradient for a linearized readout
coordinate has the form
\begin{equation}
  \left(p_0(\cdot \mid c_t) - e_{y_t}\right) \otimes \psi(c_t),
\end{equation}
where $p_0$ is the current reference distribution and $e_{y_t}$ is the one-hot
target. TSSI accumulates this quantity over the corpus:
\begin{equation}
  G_D = \sum_{t=1}^T
  \left(p_0(\cdot \mid c_t) - e_{y_t}\right) \otimes \psi(c_t).
\end{equation}
For the fixed feature map and reference distribution, $G_D$ is the sufficient
statistic used by the import rule. It is not a sufficient statistic for all
future training; it is sufficient only for the bounded question asked here:
which ternary coordinates have stable first-order CE evidence?

\subsection{Ternary Realization}

The imported artifact $\theta_I$ is a ternary sign/zero commitment. Since
$G_D[j]$ is a gradient coordinate, a positive value points toward decreasing the
parameter and a negative value points toward increasing it. The ternary import
therefore commits in the descent direction:
\begin{equation}
  \theta_I[j] =
  \begin{cases}
    -1 & \text{if } G_D[j] > \tau_j, \\
     0 & \text{if } |G_D[j]| \leq \tau_j, \\
    +1 & \text{if } G_D[j] < -\tau_j.
  \end{cases}
\end{equation}
The thresholds $\tau_j$ implement abstention. They can be set by a target
density, by calibration data, or by a coordinate-wise confidence rule. The
important property is semantic, not cosmetic: zero means ``do not claim this
coordinate yet.'' It keeps uncertain coordinates available for residual
training instead of converting weak corpus evidence into wrong signs.

\subsection{Why This Is an Import}

The imported object is smaller than a dataset and weaker than a trained model.
It contains no examples and no token-addressable memory. It also does not
contain real-valued magnitudes except through the choice of which ternary cells
are committed and which are left at zero. This is the reason the mechanism is
natural for ternary SSMs. The compiler is not trying to discover a full
continuous solution; it is trying to answer the first discrete layer of the
training problem.

\subsection{Residual Original-CE Contract}

After import, the model is still trained on the original objective. Let
$\delta$ denote the remaining trainable degrees of freedom, including
uncommitted ternary coordinates and any residual capacity. The post-import
training problem is
\begin{equation}
  \min_\delta \; \mathrm{CE}_D(\operatorname{realize}(\theta_I, \delta)).
\end{equation}
This contract is the guardrail that makes TSSI a training mechanism. If the
import writes helpful signs, residual training begins from a better ternary
artifact. If the import writes harmful signs, the original CE objective exposes
the harm. No alternative loss, teacher target, or retrieval path is introduced.

\subsection{Failure Modes}

The definition deliberately leaves room for failure. The feature map may be too
weak, the corpus may not provide stable signs, the thresholds may admit harmful
coordinates, or the imported coordinates may help a proxy but fail in live SSM
state space. These are empirical and mathematical failure modes, not
implementation details. The validation protocol in the next section is built
to attack them directly.

\section{Validation Protocol}

The validation target is not to show that an imported artifact looks good on
the same examples that created it. The target is narrower and stricter: a TSSI
artifact must improve held-out cross-entropy or held-out CE-gradient sign
alignment in a way that cannot be explained by density, sign balance, logit
scale, or example leakage.

\subsection{Identifying Checks}

The first identifying checks compare the imported artifact $\theta_I$ against a
reference start $\theta_{\mathrm{ref}}$ and density-matched random ternary
artifacts. The import is identifying only if
\begin{align}
  A_{\mathrm{sign}} &> A_{\mathrm{random}} + \mathrm{margin}, \\
  \mathrm{CE}(\theta_I) &< \mathrm{CE}(\theta_{\mathrm{ref}})
  \quad \text{on held-out data}.
\end{align}
Here $A_{\mathrm{sign}}$ measures agreement between the imported signs and
held-out CE descent signs. The second check asks whether those signs matter for
prediction, not just for a sign metric. A high alignment score without held-out
CE improvement is treated as insufficient.

\subsection{Train, Calibration, And Held-Out Separation}

The corpus stream is split into training, calibration, and held-out regions.
The training region accumulates $G_D$. The calibration region selects density,
scale, or threshold parameters. The held-out region is not used by either step.
This separation is necessary because TSSI can otherwise collapse into a
corpus-specific table that looks useful only on the evidence used to build it.
For the byte-level gates below, the main split is
300,000/100,000/100,000 tokens from \texttt{wiki\_1mb.bin}.

\subsection{Controls}

Each control removes a different easy explanation.

\paragraph{Zero or uniform start.}
This baseline tests whether the imported artifact improves the reference
cross-entropy at all.

\paragraph{Density-matched random ternary start.}
This baseline preserves sparsity and ternary occupancy while destroying corpus
placement. Passing it means the result is not simply caused by putting enough
nonzero signs into the readout or state summary.

\paragraph{Sign-shuffled start.}
This baseline preserves the empirical sign multiset while breaking the
association between signs and coordinates. Passing it means the coordinate
placement carries information.

\paragraph{Layer-summary controls.}
For live SSM states, \texttt{all}, \texttt{last}, and \texttt{mean} layer
summaries test whether the signal depends on the full recurrent state structure
or survives only as a weak pooled trace.

\subsection{Record-Disjoint Rotations}

The strongest leakage concern is that a small trace may let the compiler learn
the flattened examples rather than a reusable state-coordinate statistic.
Record-disjoint rotations address this by holding out whole trace records and
rotating the held-out pair. A result that survives every rotation is not proof
against all memorization, but it rules out the simplest explanation: one lucky
held-out suffix or direct flattened-example reuse.

\subsection{Cost Requirement}

An end-to-end training claim requires:
\begin{equation}
  R_{\mathrm{TSSI}} =
  h_{\mathrm{build}} + h_{\mathrm{realize}} + h_{\mathrm{verify}}
  + \kappa_{\mathrm{res}} \frac{H_{\mathrm{res}}}{H_0}
  \leq 0.90.
\end{equation}
This inequality is a deployment requirement, not the definition of TSSI. The
present paper focuses on whether the imported ternary artifact carries useful
cross-entropy information in model coordinates.

\subsection{Kill Criteria}

TSSI should be rejected, or at least demoted, under any of the following
conditions: held-out CE fails to beat matched random controls; sign alignment is
high but prediction does not improve; sign-shuffled controls match the import;
record-disjoint rotations collapse; live SSM coordinates lose the signal seen
in symbolic proxies; or residual CE training exposes the imported signs as
harmful. These are not secondary checks. They are part of the definition of a
useful import.

\section{Experiments}

The experiments are organized as an evidence ladder. The first gate asks
whether corpus moments can produce useful ternary signs in a simple byte-level
feature map. The second asks whether the imported artifact remains useful after
ordinary residual CE training begins. The third replaces symbolic features with
stateful recurrent features. The fourth uses live Vulkan SSM hidden states. The
last gate rotates held-out trace records to reduce the risk that the live result
is a small-trace artifact.

\subsection{Byte-Level Corpus Moment Gate}

The initial byte-level gate used \texttt{wiki\_1mb.bin} with 500,000 tokens and
a 300,000/100,000/100,000 train/calibration/held-out split. The selected
feature map was a bigram map, with density $0.05$ and scale $2.4$. The imported
artifact reached held-out CE $2.949221092857901$, compared with uniform CE
$5.545177444479562$ and random held-out CE mean $5.722178622339609$. The
weighted sign alignment was $0.9986051321029663$, versus random mean
$0.45805180817842484$.

This is the first important separation. The imported signs are not merely
nonzero ternary noise. They agree with held-out CE descent directions and reduce
held-out CE in the same gate.

\subsection{Residual Original-CE Pilot}

The byte/bigram result could still be an initialization illusion: an artifact
may look good before training, then lose its advantage once residual parameters
start moving. To test this, a residual dense readout was trained on top of fixed
imported or control logits under the same original CE target. After 1,000
residual steps, TSSI reached final CE $2.5494919515235193$. The best control
reached final CE $2.78477980855248$, giving a same-budget CE advantage of
$0.2352878570289607$. The sign-shuffled control finished at
$3.2129224560678225$ and did not reach the imported target.

This pilot matters semantically because it tests the residual contract directly.
The imported artifact does not merely lower an initial diagnostic number; it
leaves the residual learner in a better CE state than zero, random, or
sign-shuffled starts in this scope.

\subsection{Stateful Proxy and Live SSM Trace}

The next concern is that a byte bigram feature map may be too close to a direct
local statistics table. A deterministic recurrent \texttt{ssm\_state} proxy was
therefore used as a stateful feature map. With density $0.01$ and scale $1.6$,
two residual-pilot seeds preserved the advantage: seed $777$ produced
same-budget CE advantage $0.05614244344347741$, and seed $419$ produced
$0.05664628606735311$.

The stronger test used actual Vulkan SSM hidden states. The trace had shape
$10 \times 6 \times 1 \times 1024 \times 128$ and, after materialization from
\texttt{vulkan\_fwd\_h\_cache}, all $7,864,320$ hidden-state values were
nonzero. The best all-layer import used density $0.05$ and scale $3.2$,
reaching held-out CE $5.346169927602639$ versus random held-out CE mean
$5.54442637453797$. Weighted sign alignment was $0.9970957040786743$, versus
random mean $0.5062233805656433$. Under the residual pilot, the same-budget CE
advantage was $0.19843119959677402$.

The live trace is the current nucleus of the result. It moves TSSI from a
symbolic corpus statistic into the coordinate system the model actually uses.

\subsection{Record-Disjoint Anti-Memorization Gate}

The live trace is small, so the anti-memorization gate is essential. The trace
was split into records, and every two-record block was held out once under
rotations $(8,9)$, $(0,1)$, $(2,3)$, $(4,5)$, and $(6,7)$. All five all-layer
rotations passed. The CE margin versus the best density-matched random minimum
had min $0.189190727524743$, mean $0.204293772528715$, and max
$0.238266453598485$. Weighted sign alignment stayed at least
$0.9929543137550354$.

Last-layer and mean-layer summaries also passed all rotations, but with weaker
mean margins: $0.0327855562301441$ for last-layer summaries and
$0.013581927197886$ for mean-layer summaries. This supports the interpretation
that the all-layer state carries the strongest corpus-directional statistic.
The weaker summaries are useful controls, not the main result.

\section{Analysis}

The main result is not that a bigram feature map can predict bytes. That would
be too weak and too close to a conventional count table. The meaningful result
is that the same import idea survives a sequence of increasingly hostile
reinterpretations: it survives held-out CE, sign-shuffled controls, residual
training, a recurrent proxy, live SSM hidden-state coordinates, and
record-disjoint rotations.

\paragraph{The live all-layer trace is the nucleus.}
The all-layer Vulkan trace is the strongest current evidence because it places
the statistic in the actual recurrent state coordinates of the model. A
byte-level or bigram gate can be useful as an identifying probe, but it does not
by itself establish that TSSI is a model-native object. The live trace changes
the status of the idea: the corpus-derived signs are no longer only symbolic
features over token neighborhoods; they are useful over hidden states produced
by the SSM path.

\paragraph{The controls remove simple explanations.}
The density-matched random controls remove the explanation that any sparse
ternary pattern helps. The sign-shuffled controls remove the explanation that
the sign balance alone is enough. The residual pilots remove the explanation
that the import only wins before optimization begins. Record-disjoint rotations
remove the simplest small-trace leakage story. None of these checks proves the
mechanism at full scale, but together they make the positive signal hard to
dismiss as formatting, density, or one lucky split.

\paragraph{The result is still narrower than a trained model.}
TSSI should be read as a partial commitment mechanism. It does not know the
whole model. It does not decide every ternary coordinate. It does not claim that
the imported artifact is a final parameter state. Its job is to move a subset
of sign/zero decisions from iterative discovery into a corpus-statistic
compiler, while leaving the residual CE problem alive. That is exactly why the
zero state matters: the mechanism has a way to say ``not enough evidence.''

\paragraph{Why the idea is project-specific.}
The small ternary alphabet is not an incidental implementation detail. It is
what makes the import question well-posed. For a real-valued model, the corpus
would need to justify a magnitude, scale, and interaction pattern. For a
ternary SSM, the compiler can ask a sharper first question: should this
coordinate begin negative, zero, or positive under the original CE geometry?
That makes TSSI different from generic data-dependent initialization. It is not
only using data before training; it is using data to answer the discrete
commitment question that this model family actually exposes.

\section{Limitations}

The current evidence is intentionally limited. These limitations are not
decorative; they define what the paper can honestly claim.

\paragraph{Byte-level data.}
The gates use byte-level shards rather than a full BPE or unigram tokenizer
pipeline. This makes the first experiments simple and repeatable, but it leaves
open whether the same import signal holds under the tokenization and sequence
statistics used by larger language-model training runs.

\paragraph{Small live trace.}
The live Vulkan SSM evidence comes from a 10-step batch-1 trace. The trace is
valuable because it is real model state, but it is still small. Record-disjoint
rotations reduce the easiest leakage concern, yet they do not replace a larger
corpus-scale live-state validation.

\paragraph{Fixed feature maps.}
TSSI is only as good as the feature map in which corpus moments are accumulated.
A weak feature map can miss useful structure; an overly expressive feature map
can become too close to hidden retrieval. The live all-layer result is
encouraging because it uses SSM states, but the feature-map question is still
central.

\paragraph{Threshold selection.}
The current gates use density, scale, and calibration choices. These are
reasonable for identifying whether the signal exists, but a mature TSSI system
should derive or learn threshold rules that are stable across corpora, model
sizes, and training stages.

\paragraph{No downstream capability result.}
The evidence is CE-centered. It shows that corpus-derived ternary commitments
can reduce held-out CE and preserve residual CE advantage in the tested scopes.
It does not yet show downstream task improvement, long-context capability
transfer, or full training-run behavior.

\paragraph{Residual conflict remains possible.}
An imported sign can be locally useful and still interfere with later residual
learning. The current residual pilots did not expose that harm in the tested
scopes, but larger models may reveal interaction effects between imported
commitments, SSM dynamics, and later optimizer updates.

\section{Conclusion}

This paper introduced TSSI, a corpus-to-ternary-artifact mechanism for
state-space language-model training. The core idea is simple: if the model's
parameter alphabet is ternary, then a corpus can sometimes answer part of the
initial sign/zero decision problem before iterative optimization begins. The
right import object is not raw text, not a retrieval table, and not a teacher
model. It is a sufficient statistic for original cross-entropy descent in a
fixed model-coordinate feature map.

The current evidence supports the mechanism. Corpus-derived ternary
commitments beat density-matched random and sign-shuffled controls, survive
residual original-CE training, remain useful in a recurrent state proxy, and
carry signal over live Vulkan SSM hidden states. Record-disjoint rotations
further show that the strongest all-layer live-state result is not explained by
a single held-out suffix or trivial flattened-example reuse.

The next experiment is direct: build a model-native TSSI path that accumulates
the same kind of corpus moments without relying on an offline trace export,
then compare imported, zero, random, and sign-shuffled starts under residual
original-CE training at larger corpus and model scale. If that holds, TSSI
becomes more than an identifying gate. It becomes a new way to make the corpus
participate in constructing the initial ternary model state.

\bibliographystyle{plainnat}
\bibliography{references}

\end{document}
